\section{Discussion}
\label{sec:discussion}

\subsection{Relationship to Existing Frameworks}

The function-theoretic definition does not compete with the GDI, Zins'
classification, or the DIKW hierarchy---it occupies a gap each leaves open
by design. This is a deliberate scoping choice: the contribution of this
paper is structural precision at the data tier specifically, not a new
account of how data becomes information or knowledge.

\subsection{Limitations}

Two cases stress the definition's assumptions:

\paragraph{Streaming or unbounded data.}
When $I$ is not known in advance---an open-ended sensor stream, for
instance---Definition~\ref{def:collection}'s totality condition cannot be
checked against a fixed, enumerable $I$. The definition still applies at
any finite prefix or window of the stream, but the unbounded stream itself
requires treating $I$ as a fixed well-ordered set (e.g., $\mathbb{N}$)
with totality understood as ``defined for every index produced so far.''

\paragraph{Genuinely unstructured data.}
Free-text notes prior to any tagging or field extraction resist a useful
index set by design. Imposing one prematurely---treating the data as a
function from character position to character---is technically valid but
discards exactly the information a downstream NLP pipeline exists to
recover. The function-theoretic definition is correct but not yet
\emph{useful} at the byte-indexed level; its value reappears once a
downstream process imposes a more meaningful index set.

\subsection{Threats to Validity}

The corollaries are illustrated with representative, constructed examples
rather than an empirical defect-mining study. We consider the corollaries
conceptually validated---each is a direct logical consequence of
Definitions~\ref{def:collection} and~\ref{def:equality}---but their
\emph{practical prevalence} as a fraction of real-world defects is an
empirical question this paper does not measure.